Software-defined radio pushes the analog-to-digital boundary as close to the
antenna as possible and does the rest --- filtering, mixing, demodulation --- in
software.

\section{The Idea}
Instead of a fixed chain of analog filters and mixers, an SDR digitizes a wide
band and reconfigures its behavior in DSP. The same hardware can be an FM
receiver, a spectrum monitor, or a cellular base station, changed by loading new
code.

\section{Architectures}
Two dominate: \emph{direct RF sampling}, where a fast ADC digitizes the RF band
outright (using bandpass sampling); and \emph{direct-conversion}, where an
analog I/Q front end brings the band to baseband for a slower ADC. A digital
downconverter (a numerically controlled oscillator, digital mixer, and decimating
CIC/FIR filters) then selects and narrows the channel.

\section{Processing and Practicalities}
The wideband data rate off the ADC is large, so decimation reduces it to the
channel bandwidth before heavy processing on an FPGA or general-purpose
processor. SDR also enables digital predistortion, digital I/Q-imbalance and
DC-offset correction, and rapid retuning --- at the cost of ADC/DAC dynamic range
and raw data throughput.
